\subsection{ZARISKI TOPOLOGY ON Aut(g)}

PROPOSITION 8. Let V be the set of endomorphisms of the vector space g. Then \(\operatorname{Aut}(\mathfrak{g})\) is closed in V for the Zariski topology (Chap. VII, App. I).

Let \(K\) be the Killing form of \(\mathfrak{g}\) . If \(s \in \operatorname{Aut}(\mathfrak{g})\) ,

\[
  [ s x, s y ] = [ x, y ]\tag{2}
\]

\[
\mathrm{K} (s x, s y) = \mathrm{K} (x, y)\tag{3}
\]

for all \(x, y \in g\) . Conversely, let s be an element of V satisfying (2) and (3) for all \(x, y \in g\) . Then \(\operatorname{Ker}(s) = 0\) , so s is bijective and \(s \in \operatorname{Aut}(\mathfrak{g})\) . But, for all \(x, y \in g\) , the maps \(s \mapsto [sx, sy]\) and \(s \mapsto \operatorname{K}(sx, sy)\) from V to g and k are polynomial.

PROPOSITION 9. Let \(\mathfrak{h}\) be a splitting Cartan subalgebra of \(\mathfrak{g}\) .

\begin{enumerate}
\def\labelenumi{(\roman{enumi})}
\item
  The group \(f(\mathrm{T}_{\mathrm{Q}})\) is closed in \(\operatorname{Aut}(\mathfrak{g})\) in the Zariski topology.
\item
  The group \(f(q(\mathrm{T_P}))\) is dense in \(f(\mathrm{T_Q})\) in the Zariski topology.
\end{enumerate}

Assertion (i) follows from the equality \(f(\mathrm{T}_{\mathrm{Q}}) = \operatorname{Aut}(\mathfrak{g},\mathfrak{h})\cap \operatorname{Ker}\varepsilon\) (Prop. 2). Put \(m = (\mathrm{P}(\mathrm{R}):\mathrm{Q}(\mathrm{R}))\) . Let \(\mathrm{F}\) be a polynomial function on \(\mathrm{V}\) ; we assume that \(\mathrm{F}\) vanishes on the \(m\) th power of every element of \(f(\mathrm{T}_{\mathrm{Q}})\) , and show that \(\mathrm{F}|f(\mathrm{T}_{\mathrm{Q}}) = 0\) ; in view of Lemma 3, this will prove (ii).

The set \(V'\) of elements of V inducing the identity on h and leaving each \(g^{\alpha}\) stable can be identified with \(k^{R}\) . Let \(F'\) be the restriction of F to \(V' = k^{R}\) ; this is a polynomial function. We have \(f(\mathrm{T}_{\mathrm{Q}}) \subset \mathrm{V}'\) . Let \(\mathrm{B} = (\alpha_{1}, \ldots, \alpha_{l})\) be a basis of R. For all \(t = (t_{1}, \ldots, t_{l}) \in k^{*B}\) , let \(\varphi(t)\) be the homomorphism from \(\mathrm{Q}(\mathrm{R})\) to the group \(k^{*}\) that extends t. Then \(\mathrm{F}'(f(\varphi(t)))\) can be written as a finite sum

\[
\sum_ {n _ {1}, \dots , n _ {l} \in \mathbf {Z}} c _ {n _ {1}, \dots , n _ {l}} t _ {1} ^ {n _ {1}} \dots t _ {l} ^ {n _ {l}} = \mathrm{H} (t _ {1}, \dots , t _ {l}).
\]

By assumption,

\[
0 = \mathrm{H} (t _ {1} ^ {m}, \dots , t _ {l} ^ {m}) = \sum_ {n _ {1}, \dots , n _ {l} \in \mathbf {Z}} c _ {n _ {1}, \dots , n _ {l}} t _ {1} ^ {m n _ {1}} \dots t _ {l} ^ {m n _ {l}}
\]

for all \(t_{1},\ldots,t_{l}\in k^{*}\) . The \(c_{n_{1},\ldots,n_{l}}\) are thus the coefficients of a polynomial in l variables which vanishes on \(k^{*l}\) ; hence they are all zero.

PROPOSITION 10. Assume that \(\mathfrak{g}\) is splittable.

\begin{enumerate}
\def\labelenumi{(\roman{enumi})}
\item
  The group \(\operatorname{Aut}_{e}(\mathfrak{g})\) is dense in \(\operatorname{Aut}_{0}(\mathfrak{g})\) in the Zariski topology.
\item
  The groups \(\mathrm{Aut}_e(\mathfrak{g})\) and \(\mathrm{Aut}_0(\mathfrak{g})\) are connected in the Zariski topology.
\end{enumerate}

By Prop. 3, \(f(q(\mathrm{T}_{\mathrm{P}})) \subset \mathrm{Aut}_e(\mathfrak{g})\) . For all \(s \in \mathrm{Aut}_e(\mathfrak{g})\) , the closure of \(s.f(q(\mathrm{T}_{\mathrm{P}}))\) in the Zariski topology contains \(s.f(\mathrm{T}_{\mathrm{Q}})\) by Prop. 9. Hence the closure of \(\mathrm{Aut}_e(\mathfrak{g})\) contains \(\mathrm{Aut}_e(\mathfrak{g}).f(\mathrm{T}_{\mathrm{Q}}) = \mathrm{Aut}_0(\mathfrak{g})\) (Prop. 6). This proves (i).

Let \(\mathrm{Aut}_e(\mathfrak{g}) = \Omega \cup \Omega'\) be a partition of \(\mathrm{Aut}_e(\mathfrak{g})\) formed by relatively open subsets in the Zariski topology, and with \(\Omega \neq \emptyset\) . If \(\omega \in \Omega\) and if \(x\) is a nilpotent element of \(\mathfrak{g}\) , the map \(\tau : t \mapsto \omega \exp(t \operatorname{ad} x)\) from \(k\) to \(\mathrm{Aut}_e(\mathfrak{g})\) is polynomial, hence continuous in the Zariski topology; consequently, \(\tau(k)\) is connected; since \(\omega \in \tau(k)\) , we have \(\tau(k) \subset \Omega\) . Thus, \(\Omega. (\exp \operatorname{ad} kx) \subset \Omega\) , so \(\Omega. \mathrm{Aut}_e(\mathfrak{g}) \subset \Omega\) and \(\Omega = \mathrm{Aut}_e(\mathfrak{g})\) . This proves that \(\mathrm{Aut}_e(\mathfrak{g})\) is connected. It follows, by (i), that \(\mathrm{Aut}_0(\mathfrak{g})\) is connected. Q.E.D.

We shall see (§8, no. 4, Cor. of Prop. 6) that \(\mathrm{Aut}_0(\mathfrak{g})\) is closed in V in the Zariski topology, and that it is the connected component of the identity element of \(\mathrm{Aut}(\mathfrak{g})\) . On the other hand, \(\mathrm{Aut}_e(\mathfrak{g})\) is not in general closed in the Zariski topology.

* Assume that \((\mathfrak{g},\mathfrak{h})\) is split. The group \(\mathrm{Aut}_0(\mathfrak{g})\) is the group \(G(k)\) of \(k\) -points of a connected semi-simple algebraic group \(G\) with trivial centre (adjoint group). The group \(f(T_{Q})\) is equal to \(H(k)\) , where \(H\) is the Cartan subgroup of \(G\) with Lie algebra \(\mathfrak{h}\) . The inverse image \(\tilde{H}\) of \(H\) in the universal covering \(\tilde{G}\) of \(G\) (in the algebraic sense) has \(T_{P}\) as its group of \(k\) -points. The image of \(\tilde{G}(k)\) in \(G(k) = \mathrm{Aut}_0(\mathfrak{g})\) is the group \(\mathrm{Aut}_e(\mathfrak{g})\) .*

